\documentclass[../main.tex]{subfiles}
\begin{document}

%========================================================================================
% FILENAME: 10_limitations.tex (v1.0.0)
%
% §10 Limitations and Open Problems.  §10.1 six limitations (tight enumerate, no proofs);
% §10.2 two open problems via \OpenProblem (both inside the subsection so numbering reads
% 10.2.1 / 10.2.2).
%
% Silences applied; cycle only; "ABT" gone; "address" avoided in operator-disambiguation
% (it is the bound-administrator concept, not 𝔸).  limitation 6 uses \confprice (applied).
% open:attestation:leverage-timing states only the specification half (count-free, the only ratio is 1/(1-χ)).
%
% Change history lives in version control (see the versioning policy in main.tex).
%========================================================================================

\section{Limitations and Open Problems}
\label{sec:attestation:limitations}

The contract's guarantees are bounded by properties of a closed reserve pool and by what lies
outside the receipt mechanism. Six limitations and two open problems follow.

% ═══════════════════════════════════════════════════════════════════════
\subsection{Limitations}
\label{sec:limitations:limitations}
% ═══════════════════════════════════════════════════════════════════════

\begin{enumerate}[nosep]
    \item \textbf{Net economic cost is zero at $\supplyfrac{} = 1$} for all $\livesplit$ ---
        a property of closed reserve pools, not a protocol-addressable deficiency
        (\zcref{thm:costliness:conservation-identity,thm:costliness:zeta-independence}).

    \item \textbf{Maximum irrecoverable exposure during bootstrapping is
        $(1-\livesplit)\genesisreserve$} --- the time-locked-class portion whose realizable
        value depends on maturity conversion (\zcref{def:maturity:conversion}).

    \item \textbf{Portfolio-level net cost is zero.} The binding constraint is attestation
        non-transferability across addresses
        (\zcref{prop:costliness:portfolio-net-cost,lem:costliness:attestation-append-only}).

    \item \textbf{Cross-cohort deflationary asymmetry.} $\ratefloor{}$ non-decreasing
        implies late entrants pay more per unit of attestation than early entrants --- a
        reserve-model property, not addressed by the time-locked class
        (\zcref{thm:invariants:floor-non-decrease}).

    \item \textbf{External deposit incentives are exogenous.} Genuine cost discipline
        requires a competitive market (\zcref{sec:attestation:discipline}). What incentivizes external
        deposits --- agents who value the systems built on the receipt --- is outside the
        scope of the receipt mechanism.

    \item \textbf{Maturity timing is operator-discretionary.} The safeguard is the
        time-locked-class secondary-market confidence price $\confprice{k}$, which
        continuously aggregates adversarial information about the operator's intent to
        convert. Indefinite deferral suppresses $\confprice{k}$, which suppresses deposit
        inflow, which suppresses the fee revenue that motivates deferral --- a self-defeating
        strategy
        (\zcref{rem:maturity:confidence-signal,prop:discipline:anti-correlation}).
        Once announced, the maturity lead is confined to $[\matlead, \matleadcap]$
        (\zcref{postc:maturity:announcement}). Before announcement, no schedule has been
        irrevocably fixed and a valid future announcement remains possible; after
        announcement, the schedule is irrevocable. The maximum lead therefore bounds the
        post-announcement distance to scheduled conversion in cycles, while market
        discipline governs whether announcement occurs. Neither lead bound compels an
        announcement or, at this layer, guarantees wall-clock cycle progress.
\end{enumerate}

% ═══════════════════════════════════════════════════════════════════════
\subsection{Open Problems}
\label{sec:limitations:open-problems}
% ═══════════════════════════════════════════════════════════════════════

\OpenProblem{Operator Deposit Disambiguation}{operator-disambiguation}
The seigniorage results
(\zcref{thm:seigniorage:bootstrapping-bound,thm:seigniorage:no-finite-lifetime-envelope})
are stated in cumulative \emph{non-operator} deposits $D$. Distinguishing operator from non-operator
deposits requires the operator to publicly mark their own deposits. Whether this is
enforceable under this layer's public-auditability guarantee
(\zcref{postc:interface:auditability}) alone, or requires a governance mechanism whose
discharge is outside this layer's scope, is open. An alternative formulation uses total
deposits $D_{\mathrm{all}} \geq D$:
$\attest{\opidx} = \feerate\livesplit D_{\mathrm{all}}$, a looser but universally auditable
bound that attributes all deposits to non-operators.

\OpenProblem{Cross-System Burn-Timing}{leverage-timing}
% importing systems cite this as A-open:attestation:leverage-timing; the shared name is deliberate
% despite "leverage" being importing-side vocabulary (registered in the macro collision register).
Within-cycle burn timing confers an advantage bounded by $1/(1-\cycleburnfrac{})$ but not
eliminated (\zcref{def:operations:cycle-burn-fraction}). An importing system that samples
$\attest{}$ at its own boundary (\zcref{rem:interface:clocks}) and freezes a conversion
factor across its period inherits, and may amplify, this advantage; that half --- the frozen
factor itself --- is owned by the importing system's own specification.

%========================================================================================
% END OF FILE: 10_limitations.tex (v1.0.0)
%========================================================================================

\end{document}
